\chapter{Gerätelehre --- \gAasRsN{RS~IX}}

\emph{Maintainer: Roman Koch}

\minitoc

\section{Medizinproduktegesetz - MPG}

\gMarried{
Das Medizinproduktegesetz regelt die Funktionstüchtigkeit,
Sicherheit, Inbetriebnahme, Instandhaltung, den Betrieb und
die Anwendung, Überwachung, Sterilisation, Desinfektion und
Reinigung von Medizinprodukten. Als Medizinprodukte gelten
alle Instrumente, Apparate, Vorrichtungen, Stoffe oder
andere Gegenstände (inkl. Software) die vom Hersteller zur
Anwendung am Menschen bestimmt sind. Es gibt über 400.000
verschiedene Produkte. Doch eines haben alle gemeinsam: die
\gE{CE-Kennzeichnung}.
}{
\centering \includegraphics[width=0.5\linewidth]{./images/geraete/Ce-logo.jpg}
\captionof{figure}{Das CE-Kennzeichen.}
}


\paragraph{Einteilung in Risikoklassen}
\begin{center}
\begin{tabulary}{\linewidth}{lLL}\toprule
\gEs{Klasse I}
&
Kein oder geringes Anwendungsrisiko
&
z.B: Verbandsmaterial, Gehhilfe, Rollstuhl, Tragsessel, Sehhilfe, Stethoskop, Trage, Beatmungsbeutel, Fieberthermometer, Druckminderer
\\\midrule
\gEs{Klasse IIa}
&
Anwendungsrisiko ohne entsprechender Ausbildung bzw. Unterweisung
&
z.B: aktive Diagnosegeräte, Absaugeinheit und -katheter, EKG-Kabel, Zuckerstreifen, Trachealtuben, Dauerkatheter
\\\midrule
\gEs{Klasse IIb}
&
Erhöhtes Anwendungsrisiko ohne entsprechender Ausbildung bzw. Unterweisung
&
z.B: Beatmungsgerät, Defibrillator, Pulsoxymeter, Babyinkubator, Röntgengeräte
\\\midrule
\gEs{Klasse III}
&
Besonders hohes Anwendungsrisiko ohne entsprechende Ausbildung bzw. Unterweisung
&
z.B: Perfusor, Insulinpumpe, Herzkatheter oder -klappen, Medizinprodukte mit Arzneimittelkomponenten
\\\bottomrule
\end{tabulary}
\end{center}

\paragraph{Pflichten des Sanitäters gemäß MPG}

Sanitäter sind lt. MPG verpflichtet, die Medizinprodukte
richtig und sicher einzusetzen bzw. anzuwenden, so wie es
ihnen bei ihrer Ausbildung bzw. Unterweisung gelehrt wurde.
Ausserdem sind sie verpflichtet die Geräte vor dem Gebrauch
mittels eines Gerätecheck auf ihre Funktion zu überprüfen.

\section{Blutdruck- und Herzfrequenzmessung}

\label{topic:rr-messung}

% Grundlagen
% 
% Fallstricke
% 
% Die Blutdruckmanschette soll nach M\"oglichkeit nicht auf der Seite
% 
% \begin{itemize}
%     \item eines Shuntarmes bei Dialysepatienten
%     \begin{itemize}
%         \item Ein Dialyse-Shunt ist ein künstlicher Kurzschluss zwischen einer
%         Armarterie und einer Armvene. Dadurch wird die Vene erweitert und sie
%         pulsiert, mann kann viel mehr Durchfluss für die Dialyse erreichen.
%         Die Nahtstelle zwischen Vene und Arterie ist jedoch sehr sensibel, wenn
%         gestaut wird kann sie aufplatzen 
%         \item ${\rightarrow}$ Gefahr einer massiven Shuntblutung
%     \end{itemize}
%     \item eines Arm\"odems
%     \item einer Brust-OP mit Lymphknoten-Entfernung (\"Odemgefahr)
%     \item einer Halbseitenschwäche oder -lähmung
% \end{itemize}
% 
% angelegt werden. 
% 
% \gCave{Shunt-Arme sind tabu für die Blutdruckmessung!}
% 
% Bitten Sie im Praktikum am KTW Dialysepatienten h\"oflich Ihnen ihren
% Shunt-Arm zu zeigen!
% 
% Ma{\ss}nahmen
% 
% \begin{enumerate}
%     \item Evt. Funktion des Blutdruckmessgerätes und des Stethoskops prüfen
%     \begin{itemize}
%         \item Kontrolle der Dichte des Ventils, der Verbindungsschläuche und
%         der Manschette (Ventil schlie{\ss}en, Manschette etwas aufpumpen und
%         Zusammenpressen, Manschette ganz entlüften)
%         \item Kontrolle der Gängigkeit des Ventils (es soll sich leicht auf
%         und zuschrauben lassen)
%         \item Kontrolle des Manometers auf korrekte Eichung: Die Nadel muss sich
%         im Ruhezustand genau bei 0\,mmHg befinden
%         \item Ohrst\"opsel des Stethoskops zeigen schräg nach vorne
%         \item Trichterfunktion durch Berühren der Membran prüfen
%     \end{itemize}
%     \item Patientenlagerung
%     \begin{itemize}
%         \item Arm freimachen und entspannt auf einen Tisch legen in leichter
%         Flexion und Supination
%         \item Oberarm auf Herzh\"ohe
%     \end{itemize}
%     \item Auswahl der geeigneten Manschettengr\"o{\ss}e
%     \begin{itemize}
%         \item Manschette laut Herstellerangaben wählen z.B Standardmanschette:
%         Oberarmumfang max. 31 cm
%         \item Messwertverfälschung vermeiden! Der aufblasbare Teil der
%         Manschette soll 80\% des Oberarmumfangs umfassen, die Breite der
%         Manschette mindestens 1/3 des Oberarmumfangs betragen
%         \item Zu kleine Manschetten führen zum Messen zu hoher Werte, zu weite
%         Manschetten k\"onnen nicht am Arm fixiert werden
%     \end{itemize}
%     \item Handhabung des Stethoskops
%     \item Bestimmen der Anlegestelle für Manschette und der
%     Auskultationsstelle durch Palpation
%     \begin{itemize}
%         \item Arteria brachialis medial in der Ellenbeuge aufsuchen
%     \end{itemize}
%     
%     \item Anlegen der Manschette
%     \begin{itemize}
%         \item Manschette muss auf unbekleideter Haut aufliegen
%         \item mä{\ss}ig straff
%         \item Abstand zur Ellenbeuge 2,5 cm
%         \item aufblasbaren Teil über A. brachialis zentrieren
%         \item Schläuche nicht im Weg,
%         \item CAVE: Halbseitenschwäche/-lähmung, Shuntarm bei
%         Dialysepatienten, Arm\"odem
%     \end{itemize}
%     \item Puls fühlen beim Aufpumpen
%     \begin{itemize}
%         \item Schnelles Aufpumpen
%         \item Tasten des Radialispulses mit Finger 2-4
%         \item Aufpumpen bis 30 mmHg über jenen Druck, ab dem kein Radialispuls
%         mehr getastet wird
%     \end{itemize}
%     \item Messvorgang
%     \begin{itemize}
%         \item Kopf des Stethoskops auf a. brachialis legen
%         \item Druck langsam ablassen 2 mmHg/sec
%         \item Systolischer Wert: Beginn der turbulenten Korotkoffschen
%         Str\"omungsgeräusche
%         \item diastolischer Wert: Verschwinden der Korotkoffschen Geräusche
%         \item Wenn Geräusche überhaupt nicht verschwinden (z.B. bei Kindern,
%         Schwangeren etc.), so gilt als diastolischer Wert der Messpunkt, an dem
%         die Geräusche deutlich leiser werden
%         \item Wenn Geräusche verschwunden sind, noch weitere 10mmHg langsam
%         ablassen -- wenn dann keine Geräusche mehr h\"orbar sind, kann die
%         Manschette zügig entleert werden
%         \item CAVE: auskultatorische Lücke
%     \end{itemize}
%     \item Protokollierung der Blutdruckwerte
% \end{enumerate}



\subsection{Messung der Herzfrequenz -- HF}

Als Herzfrequenz (\gT{HF}) bezeichnet man die Anzahl der
Schläge des Herzens pro Minute. Sie wird mittels Fühlen des Pulses
ermittelt. Dieser wird idealerweise an der
\gEs{Handgelenkarterie} (A. \gT{radialis}) bzw. bei nicht
ansprechbaren oder schockierten Patienten an der
\gEs{Halsschlagader} (A. \gT{carotis}) ertastet. Bei
Neugeborenen und Säuglingen kann man auch an der
\gEs{Oberarmarterie} (A. \gT{brachialis}) fühlen. Gemessen
wird immer mit zwei bis drei Fingern. Nie jedoch mit dem Daumen, da hier meistens nur der eigene Puls gefühlt wird und nicht der des
Patienten. Wenn an der Halsschlagader (A. carotis) die
Herzfrequenz ermittelt wird, darf nie auf beiden Seiten
gleichzeitig gemessen werden, da sonst die Gefahr besteht,
dass es zu einer Unterversorgung des Gehirn durch Blut kommen
kann.

\paragraph{Durchführung}

Es wird an der Handgelenkarterie (A. radialis) der Puls mit
zwei bis drei Fingern ertastet. Nun werden 15 Sekunden lang
die Schläge gezählt. Multipliziert man den Messwert mit vier, erhält man  die Herzfrequenz (HF) des Patienten. Der Normalwert liegt beim Erwachsenen zwischen 60 und 100 Schlägen pro Minute.
Dokumentiert wird sie mit einem HF (z.B.: HF=70)

\subsection{Messung des Blutdrucks -- RR}

\gMarried{Der Blutdruck ist jener Druck, welcher in unseren Blutgefäßen
herrscht. Wenn wir vom "`Blutdruck"' sprechen beziehen wir
uns auf den Blutdruck in den
Arterien\footnote{\emph{Arterie:} Vom Herzen wegführendes
Blutgefäß.}. Er ist in Kombination mit der Herzfrequenz
einer der wichtigsten Messwerte im
Rettungsdienst.\vspace{1em}

Der Blutdruck kann mittels elektronischer
Blutdruckmessgeräte ermittelt werden. Diese sind aber im
Rettungsdienst nicht in Verwendung, da es immer wieder zu
Gerätefehler kommen kann. Außerdem können die Batterien leer
werden. Solche Geräte haben meistens die Patienten für die
tägliche Überwachung ihres Blutdrucks. Leider bedienen aber
viele die Geräte falsch und fürchten sich dann
lebensbedrohlich krank zu sein.

Im Rettungsdienst wird die Blutdruckmessung nach Riva-Rocci
durchgeführt. Benannt nach dem Erfinder Scipione
Riva-Rocci. Dabei wird eine aufblasbare Manschette mit
angeschlossenem Dosenmanometer und Pumpbällchen benötigt.
Dieser Teil wird in der Praxis Riva-Rocci genannt. Im
Idealfall wird auch noch ein Stethoskop verwendet. }
 {\begin{center}
    \includegraphics[width=4cm]{./images/noauth/GER-rr-messgeraet.jpg}
 
 
  \includegraphics[width=4cm]{./images/noauth/GER-rr-elektrisch-oberarm.png}
  
   \includegraphics[width=4cm]{./images/noauth/GER-rr-elektrisch-unterarm.jpg}
  \end{center}
 }

\paragraph{Messung nach Riva-Rocci}

Bei der Blutdruckmessung gibt es zwei Möglichkeiten zu
einem Wert zu kommen. Die genauere ist die
\gEs{auskultatorische} Methode (durch hören). Die zweite
Möglichkeit ist die \gEs{palpatorische} Methode (durch
tasten). Im Idealfall werden zwei Werte ermittelt. Der
\gEs{systolische Wert} (obere Wert) ist jener Druck im
Blutgefäßsystem während der Herzaktionsphase, in der sich die
Herzkammern anspannen und zusammenziehen (Auswurfphase). Der
\gEs{diastolische Wert} (untere Wert) ist jener Druck im
Blutgefäßsystem während der Herzaktionsphase, in der sich die
Herzkammern entspannen und mit Blut füllen
(Entspannungsphase). Der diastolische Wert kann nur
mittels der auskultatorischen Methode ermittelt werden. Die
\gEs{Einheit}, in der der Blutdruck angegeben wird, ist
\gT{Millimeter Quecksilbersäule}
(\gT{mmHg}).\index{Milimeter Quecksilbersäule}\index{mmHg}

\minisec{Durchführung}

\paragraph{Auskultatorische Methode}

Die Manschette sollte faltenfrei, nicht über Kleidung, in der
Mitte des Oberarms angelegt werden. Nun wird der Puls an der
Handgelenkarterie (a.radialis) ertastet. Die Manschette wird
solange aufgepumpt, bis der Puls nicht mehr getastet werden
kann. Nun noch 20\,mmHg weiter aufpumpen. Das Stethoskop
wird in der Ellenbeuge auf die Oberarmarterie (a.brachialis)
angelegt. Die Luft langsam aus der Manschette ablassen.
Sobald man ein Pochen hört (Korotkowsche Geräusche) muss man
den Messwert am Dosenmanometer ablesen. Dies ist der
\gEs{systolische} Wert. Die Luft wird weiter abgelassen, und
wenn das Pochen verschwindet,  muss der zweite Wert abgelesen
werden. Dies ist der \gEs{diastolische} Wert. Die Werte
werden immer auf die nächste Fünferstelle abgerundet (z.B.:
123/83\,mmHg wird abgerundet auf 120/80\,mmHg). Dokumentiert
wird der Blutdruck mit RR (z.B.: RR=125/75\,mmHg). Der
Normalwert beim Erwachsenen liegt bei RR 120/80\,mmHg.


\begin{figure}[H]
    \begin{center}
   
   
    \includegraphics[width=10cm]{./images/noauth/GER-rr-geraeusche.png}

    \captionof{figure}{\gAuthNotesPicLegalNotOk{GER-rr-geraeusche.png}{}}
    \end{center}
\end{figure}




\paragraph{Palpatorische Methode}
Die Manschette sollte faltenfrei, nicht über Kleidung, in der
Mitte des Oberarms angelegt werden. Nun wird der Puls an der
Handgelenkarterie (A. radialis) ertastet. Die Manschette
wird solange aufgepumpt bis der Puls nicht mehr getastet werden
kann. Nun noch 20\,mmHg weiter aufpumpen. Die Finger bleiben
auf der Handgelenkarterie (A. radialis). Die Luft langsam
aus der Manschette ablassen. Sobald man wieder einen Puls fühlt,
muss man den Messwert am Dosenmanometer ablesen. Dies ist der
\gEs{systolische} Wert. Der \gEs{diastolische} Wert kann mit
dieser Methode nicht ermittelt werden. Für die Notfallmedizin
ist in der Regel der systolische Wert der wichtigere.

\minisec{Was ist noch zu beachten?}

Die Blutdruckmanschette soll nach Möglichkeit nicht auf der Seite
\begin{itemize}
	\item eines Shuntarmes bei Dialysepatienten
	\begin{itemize}
		\item Ein Dialyse-Shunt ist ein künstlicher Kurzschluss
		zwischen einer Armarterie und einer Armvene. Dadurch wird
		die Vene erweitert und sie pulsiert. Mann kann viel mehr
		Durchfluss für die Dialyse erreichen. Die Nahtstelle
		zwischen Vene und Arterie ist jedoch sehr sensibel, wenn
		gestaut wird kann sie aufplatzen $\rightarrow$ Gefahr
		einer massiven Shuntblutung.
	\end{itemize}
		\item eines Armödems
		\item einer Halbseitenschwäche oder -lähmung 
\end{itemize}
angelegt werden.

\gCave{Shunt-Arme sind tabu für die Blutdruckmessung!}




\section{$O_2$ --
Sauerstoff}\label{topic:sauerstoff}\label{topic:o2}\index{Sauerstoff}\index{$O_2$}

\gDefinition{Sauerstoff (chem. Symbol $O_2$) ist ein
lebenswichtiges Gas welches zu 21\% in unserer Atemluft vorkommt und au{\ss}erdem ein essentielles Notfall{\textquotedblright}medikament{\textquotedblleft}. Die
frühzeitige Verabreichung ist wichtig für die Prognose des
Patienten!}

Sauerstoff darf und muss bei Bedarf vom RS verabreicht werden.

% \begin{center}
% \includegraphics[width=4.522cm,height=5.628cm]{./images/rippenspreitzer-02-fett.png}
% \captionof{figure}{NEIN! \gSourcePicture{Rippenspreitzer.de}}
% 
% \end{center}

\subsection{Sauerstoff -- technische Grundlagen}
Sauerstoff wird in Druckflaschen  gelagert. Laut EU-Norm sind diese
Flaschen wei{\ss} und mit einem {\quotedblbase}N{\textquotedblleft}
gekennzeichnet. Das {\quotedblbase}N{\textquotedblleft} steht für
{\quotedblbase}neue Kennzeichnung{\textquotedblleft}, da früher die
Flaschen blau gekennzeichnet wurden. \footnote{{\textquotedblright}N{\textquotedblleft} bedeutet in diesem Fall nicht Stickstoff!}

Die Druckflaschen haben unterschiedliche Füllgr\"o{\ss}en (Volumina),
gängige Flaschengr\"o{\ss}en sind 0,5, 1, 1,5, 2, 5 und
10\,l. Der Sauerstoff wird unter Druck gelagert, wodurch die
Menge des zur Verfügung stehenden Sauerstoffs um ein Vielfaches
gesteigert wird. Der Druck einer gefüllten Flasche beträgt
normalerweise ca. 200\,bar. 

Jede Druckflasche verfügt über ein \gT{Hauptventil} und
ein Standardgewinde. An diesem Gewinde ist entweder ein
Druckschlauch oder ein \gT{Druckminderer} angeschlossen. Der 
Druckminderer hat ein \gE{Abgabeventil}, welches den
\gE{Druck drosselt} und eine Abgabe an den Patienten
erm\"oglicht. Je nach Bauart kann der Druckminderer über
ein \gT{Regelventil}, ein \gT{Flowmeter}\index{Flowmeter} und ein
\gT{Manometer}\index{Manometer} verfügen. Mit dem \gE{Regelventil} kann man
die Durchflu{\ss}- bzw. Abgaberate ein\-stellen, das
\gE{Flowmeter} (Durchflu{\ss}anzeige) zeigt die aktuelle
Durchflu{\ss}rate in Liter pro Minute (\gLMin) an. Das
\gE{Manometer} zeigt den Druck in der Sauerstofflasche an
wenn das Haupt\-ventil ge\"offnet ist.

\subsection{Heiteres Rechnen mit Sauerstoff}

Mit Sauerstoff lässt sich gut rechnen ...

Die insgesamt verfügbare Menge wird errechnet indem man
die Füllgr\"o{\ss}e der Flasche (in~l) mit dem
Flaschendruck (in bar) multipliziert.

\begin{equation}
\mbox{Verfügbares Volumen in l} = {\begin{array}{c}
\mbox{Volumen der Flasche in l}\\
\mbox{(Füllgröße)} \end{array}}
\times 
\begin{array}{c}{\mbox{Druck in bar}}\\
\mbox{(Flaschendruck)}
\end{array}
\end{equation}


%oder kurz: 

%\begin{equation}
%l = {Volumen} \times {Druck}
%\end{equation}

Wenn man wissen m\"ochte wie lange die vorhandene Sauerstoffmenge bei
einer bestimmten Abgaberate reicht, dividiert man einfach
die gesamt verfügbare Menge durch die Abgaberate. Achtung:
Es ist eine \gEs{ausreichende Reserve} einzuplanen, und es
ist zu Bedenken, dass in der Flasche ein \gEs{Restdruck}
verbleiben muss!

\begin{equation}
\mbox{Zeit} = \frac{\mbox{Verfügbare
Menge}}{\mbox{Abgaberate}}
\end{equation}

%oder kurz: 

%\begin{equation}
%min = \frac{l}{l / min}
%\end{equation}


\minisec{Beispiel}

\begin{quotation}
Flasche 10\,l mit 150\,bar: 1500\,l  $O_2$

Patient ben\"otigt 5\,\unitfrac{l}{min} $O_2$ : 

Die Füllung reicht für 300 Minuten, abzüglich eines
Sicherheitsvolumens.
\end{quotation}

\subsection{Umgang mit Sauerstoff und Druckflaschen}

Die Sauerstoffflaschen sind pfleglich zu behandeln!

\gCave{Sauerstoff ist zusammen mit Fett selbstentzündlich!
Die Armaturen und Gewinde sind fettfrei zu halten! Ventile dürfen nicht mit fetten oder öligen Händen
bedient werden! } 


\gCave{Das Hantieren mit Feuer und offenem
Licht ist in der Nähe von Sauerstoffflaschen verboten!} 

Der gro{\ss}e Druck birgt weitere Gefahren: Sollte ein Ventil abbrechen
kann die Flasche eine enorme Kraft entwickeln und zu einer Rakete
werden. Die Kraft einer mit 200\,bar gefüllten Flasche reicht aus um
Mauern zu druchbrechen!  

\gCave{Achtung! Vorsicht beim Hantieren -- Zerknall- und 
Explosionsgefahr!}



\subsection{Verabreichungsarten}

Grundsätzlich unterscheidet man zwischen Berieselung und
Beatmung: 

\subsubsection{Berieselung}
Nur bei vorhandener Spontanatmung!

% zwei kurze, in die Nasen\"offnungen  eingeführte Kunst\-stoff\-stutzen
% 
% bei ca. 4\,l O2 /min bis 40\% $O_2$  in Atemgas  m\"oglich
% 
% 
% Bei behinderter Nasenatmung kann die  Sauerstoffbrille im Mundwinkel
% plaziert werden. 

\begin{center}
\begin{tabulary}{\linewidth}[H]{CLLL}
\toprule
Was
 &
Wann
 &
Vorteile
&
Nachteile
\\\toprule 
\gT{Sauerstoff\-brille}
 &
Spontanatmung

$O_2$-Flow {\textless} 5\gLMin
 &
gute Toleranz

fehlendes Engegefühl

Sprechen\&Husten m\"oglich
 &
ungenaue Dosierung

Austrocknung  der Schleimhäute
\\\midrule
\gT{Sauerstoff\-maske}
 &
Spontanatmung

$O_2$-Flow ab 5\gLMin
 &
h\"ohere $O_2$-Konzentration
 &
$CO_2$-Rückatmung  wenn $O_2$-Durchflu{\ss}
{\textless} 5\gLMin!) 

für Pateinten evt. unangenehm bzw. gewöhnungsbedürftig
\\\midrule
\gT{Sauerstoff\-maske mit Reservoir}
 &
Spontanatmung

$O_2$-Flow ab 5\gLMin
 &
Noch h\"ohere $O_2$-Konzentration
 &
Wie oben.

Reservoir muss zuerst händisch befüllt werden.
\\\bottomrule
\end{tabulary}
\end{center}

\gDias
{\includegraphics[width=2.5cm]{./images/o2-brille.png}\captionof{figure}{Sauerstoffbrille}}
{\includegraphics[width=2.5cm]{./images/o2-maske.png}\captionof{figure}{Sauerstoffmaske}}
{\includegraphics[width=2.5cm]{./images/o2-maske-reservoir.png}
\captionof{figure}{Sauerstoffmaske mit Reservoir}}

\subsubsection{Beatmung}

Bei fehlender Spontanatmung!

Entweder mittels Beatmungsbeutel (s.
\ref{topic:beatmungsbeutel}) oder Beatmungsgerät. 

\begin{center}
\includegraphics[width=7.189cm,height=4.94cm]{./images/beatmungsgeraet-medumat.png}
\captionof{figure}{Beatmungsgerät \emph{Medumat
Standard\textregistered}. \gSourcePicture{ASB
Floridsdorf-Donaustadt}}
\end{center}



Mit dem Beatmungsgerät \emph{"`Medumat"'} der Fa. Weinmann
kann ein Patient sowohl beatmet als auch berieselt werden
(sofern ein Berieselungsmodul vorhanden ist). 

Dafür gibt es zwei Bedienfelder: Das eine zum Berieseln,
welches vom Sanitäter eingestellt und verwendet werden
kann; das zweite darf nur auf Anweisung eines Notarztes
verwendet werden und wird für das Beatmen benötigt. 


\includegraphics[width=\linewidth]{./images/geraete/ger-medumat-standard-bedienfeld.jpg}

% \gCave{JEDE BEATMUNG SOLLTE MIT RESERVOIRBEUTEL UND 10-15\gLMin
% $O_2$ DURCHGEF\"UHRT WERDEN!}
% 
% Beatmungsfrequenz beim Erwachsenen: 12-15\,x\,/\,min (Eigenfrequenz)
% 
% Achtung! Zuviel ist schlecht! Nicht zu schnell beatmen ${\rightarrow}$
% Hyperventilationsgefahr!
% 
% \subsection{Sauerstofftherapie}
% Prinzipiell sollte sich die gewählte Therapie an der Diagnose und am
% Schweregrad orientieren!
% 
% Bedenke: 
% 
% \begin{itemize}
% \item Diagnose!
% 
% \item Schockbekämpfung
% 
% \item Ischämie (bei Akutem Koronoarsyndrom, Insult, ..)
% 
% \item Atemst\"orungen (Fremdk\"orper, Atemnot, {\dots})
% 
% \end{itemize}

\section{Beatmungsbeutel}
\label{topic:beatmungsbeutel}

\gMarried{
Im Rettungsdienst sollte jede Beatmung mit einem
Beatmungsbeutel oder einem Beatmungsgerät durchgeführt
werden. Wenn an dem Beatmungsbeutel dann noch ein
Verbindungsschlauch an eine Sauerstoffflasche mit einem Flow
von 15\gLMin angeschlossen wird und zusätzlich auch noch ein
Reservoir angesteckt ist,  wird mit der größt möglichen
Reserve an Sauerstoff beatmet.
}{
\begin{center}
\includegraphics[width=\linewidth]{./images/geraete/ger-ambu.jpg}
\captionof{figure}{Beatmungsbeutel mit Reservoir,
Bakterienfilter und Masken unterschiedlicher Größe.}
\end{center}
}



\begin{center}
\begin{tabulary}{\linewidth}{Lr}\toprule
\gTabHeading{Mögliche Sauerstoffkonzentrationen bei der Beatmung} 
\\\midrule Mund-zu-Mund Beatmung
&
17\,\%
\\\addlinespace
Beutel-Masken-Beatmung
&
21\,\%
\\\addlinespace
Beutel-Masken-Beatmung mit angeschl. $O_2$ (15\gLMin)
&
40 - 70\,\%
\\\addlinespace
Beutel-Masken-Beatmung mit angeschl. $O_2$ (15\gLMin) und
Reservoir &
fast 100\,\%
\\\bottomrule
\end{tabulary}
\end{center}

\gMarried{
Ein Beatmungsbeutel besteht im Idealfall aus folgenden Teilen:
\begin{itemize}
	\item Beatmungsmaske 
	\item Bakterienfilter
	\item Ein-/Ausatemventil mit angeschlossenem PEEP-Ventil
	\item Beatmungsbeutel
	\item Reservoir mit einem Verbindungsschlauch mit einer Sauerstoffflasche verbunden
\end{itemize}
}{
\includegraphics[width=\linewidth]{./images/geraete/ambu-komplett-zerlegt.jpg}
}


\subsection{PEEP - Ventil}

\gDefinition{PEEP = Positive EndExspiratory Pressure =
positiver end-exspiratorischer Druck}

\gMarried{
Das PEEP-Ventil wird beim Ein-/Ausatemventil des
Beatmungsbeutels angeschlossen. Es hält einen positiven Druck
in der Lunge aufrecht. Nachdem die Verwendung des
PEEP-Ventils eine notärztliche Maßnahme ist, muss es bei der
Anwendung durch einen Sanitäter immer auf Null stehen.
Mögliche Einstellungen für das PEEP-Ventil sind 5 bzw. 10 cm
Wassersäule.
}{
\includegraphics[width=0.6\linewidth]{./images/geraete/peep-ventil.jpg}
}

\gCave{Ein PEEP-Ventil ist kein Einmalprodukt!}


\section{Absauggeräte}

Blut oder Schleim können von Sanitätern mittels Absaugpumpen
und dem entsprechenden Absaugkatheter aus dem
Mund-Rachen-Raum abgesaugt werden. Es gibt unterschiedliche
Modelle von Absaugpumpen. Grundsätzlich wird nach der
Betriebsart unterschieden zwischen \gEs{elektrischer
Absaugeinheit}, \gEs{manueller Absaugpumpe} und als
Sonderfall nur für eine bestimmte Patientengruppe geeignet
der \gEs{Oro-Sauger} oder Schleimabsauger.

\begin{center}
\begin{tabulary}{\linewidth}{Lp{4cm}}\toprule
\gTabHeading{Elektrische Absaugeinheit}

älteres Model von der Fa. Lerdal (Suction Unit)

neueres Model von der Fa. Weimann (Accuvac)
&
\vspace{0.1pt}
\includegraphics[width=\linewidth]{./images/geraete/lsu-alt.jpg}

\includegraphics[width=\linewidth]{./images/geraete/accuvac.jpg}
\\\midrule
\gTabHeading{Manuelle Absaugpumpen}

Handabsaugpumpe

Fußabsaugpumpe
&
\vspace{0.1pt}
\includegraphics[width=\linewidth]{./images/geraete/handabsaugung.jpg}

\includegraphics[width=\linewidth]{./images/geraete/fussabsaugung1.jpg}
\\\midrule
\gTabHeading{Oro-Sauger}

Schleimabsauger für Neugeborene
&
\vspace{0.1pt}
\includegraphics[width=\linewidth]{./images/geraete/orosauger.jpg}
\\\midrule
\gTabHeading{Absaugkatheter}

zum Einführen in Mund oder Nase
&
\vspace{0.1pt}
 \includegraphics[width=4cm]{./images/geraete/absaugkatheter.jpg} \\\bottomrule
\end{tabulary}
\end{center}

Die Absaugkatheter sind steril verpackt und sollten auch
dementsprechend behandelt werden. Die farbige Kennung gibt
den Durchmesser des Absaugkatheters an. Wenn eine
\gEs{Absaugbereitschaft} hergestellt werden soll, wird nur
der farbige Teil des Katheters ausgepackt und an den
Konnektor bzw.\gEs{Fingertipp} der Absaugpumpe angeschlossen.
Der restliche Katheter bleibt noch verpackt. Damit der
Absaugkatheter nicht zu tief eingeführt wird, kann die
einzuführende Katheterlänge durch Abmessen der Distanz
zwischen Mundwinkel und Ohrläppchen bzw. zwischen Nasenspitze
und Ohrläppchen bestimmt werden.

\gCave{Es darf nur auf Sicht abgesaugt werden!}

\begin{figure}[H]
    \begin{center}

    \includegraphics[width=10cm]{./images/noauth/GER-absaugen.png}
 
    \captionof{figure}{\gAuthNotesPicLegalNotOk{GER-absaugen.png}{}}
    \end{center}
\end{figure}


\section{Transporttechniken bzw. -geräte}

Als Sanitäter hat man während des gesamten Transports die
Verantwortung für den Patienten. Dies gilt auch für gehfähige
Patienten. Sollte der Patient auf Grund seines hohen Alters
oder seiner Kreislaufsituation nicht selbstständig gehen
können, muss er entweder \gEs{sitzend} oder \gEs{liegend}
transportiert werden.

\subsection{Techniken bzw. Geräte für den sitzenden Transport} 

\paragraph{Tragen mit dem Tragering}

Funktioniert nur bei bewusstseinsklaren und kooperativen
Patienten. Ausserdem sollte das Stiegenhaus oder die zurück
zulegende Strecke breit genug sein. Eignet sich nur für kurze
Distanzen.

\paragraph{Umgang mit dem Tragsessel}~


\gMarried{
Der Patient ist immer anzuschnallen.Beim Transport über
Stiegen hat der Patient immer talwärts zu blicken.
}{
\begin{center}
    \includegraphics[width=4cm]{./images/noauth/GER-tragsessel.jpg}
    \captionof{figure}{\gAuthNotesPicLegalNotOk{GER-tragsessel.jpg}{}Tragsessel}
\end{center}
}

\paragraph{Umgang mit einem Rollstuhl}

Siehe Abb. \ref{fig:rollstuhl-umgang}.

~

\subsection{Techniken bzw. Geräte für den liegenden Transport}

\paragraph{Tragetuch}

Das Tragetuch (ESE-Tuch) eignet sich ideal für den Transport
von liegenden Patienten wo der Einsatz einer Krankentrage
eventuell aus Platzgründen nicht möglich ist. Es wird auch
gerne zum Umlagern von Patienten oder zum Retten aus einer
Gefahrenzone verwendet. Der Patient wird immer mit den Füßen
in Transportrichtung getragen.

% \begin{figure}[H]
%     \begin{center}
%  TODO Foto Tragetuch
%     \includegraphics[width=10cm]{./images/noauth/GER-tragetuch-tragen.png}
% 
%     \captionof{figure}{\gAuthNotesPicLegalNotOk{GER-tragetuch-tragen.png}{}}
%     \end{center}
% \end{figure}

%\gTempPicMissing{EINFÜGEN BILDER TRAGETUCH VON PP}

\paragraph{Einheitskrankentrage mit Fahrgestell}~

\gMarried{
Der Patient ist immer anzuschnallen. Beim Transport über
Stiegen hat der Patient immer talwärts zu blicken.
}{
    \begin{center}

   
    \includegraphics[width=7cm]{./images/noauth/GER-fahrtrage-stollenwerk.jpg}
   
   
    \captionof{figure}{\gAuthNotesPicLegalNotOk{GER-fahrtrage-stollenwerk.jpg}{}Krankentrage mit Fahrgestell.}
    \end{center}
}


\paragraph{Umlagern von Patienten}

    \begin{center}

    \includegraphics[width=10cm]{./images/noauth/GER-umlagern.png} 
    \captionof{figure}{\gAuthNotesPicLegalNotOk{GER-umlagern.png}{}\label{fig:rollstuhl-umgang}\label{fig:umlagern}}
    \end{center}


\section{Geräte zum Retten von Patienten}

\subsection{Schaufeltrage}
\index{Schaufeltrage!Beschreibung}

\gMarried{
Die Schaufeltrage wird primär für das Umlagern von liegenden
Patienten mit dem Verdacht auf eine Verletzung der
Wirbelsäule verwendet. Das Anwendungsgebiet ist jedoch sehr
vielfältig. Es reicht vom sicheren Transport im unwegsamen
Gelände bis zum Retten aus einer LKW-Kabine.

\vspace{3em}
    \captionof{figure}{\gAuthNotesPicLegalNotOk{GER-schaufeltrage.jpg}{}Eine Schaufeltrage}
}{
    \begin{center}
    \includegraphics[width=4cm]{./images/noauth/GER-schaufeltrage.jpg}   
    \end{center}
}

\subsection{Spineboard}
\label{topic:spineboard-beschreibung}\index{Spineboard!Beschreibung|gZ}

Das Spineboard ist im Prinzip eine flaches Brett auf dem
der Patient ähnlich dem Schufeltrage/Vakuummatratze-Duett
fixiert werden kann. Diese Technik wurde ursprünglich vor
allem im Anglo-amerikanischen Raum eingesetzt und findet
auch hierzulande immer größere Verbreitung.

\ASBFLO{Das Spineboard gehört derzeit nicht zur
Standardausstattung und wird nicht Routinemäßig verwendet.}

\gTakeNotesMedium

\subsection{Rettungskorsett}
\index{Rettungskorsett!Beschreibung}

\gMarried{
Das Rettungskorsett wird primär für das Umlagern von
sitzenden Patienten mit dem Verdacht auf eine Verletzung der
Wirbelsäule verwendet. Es kann aber auch für das Retten
aller sitzenden Patienten aus einer Gefahrenzone verwendet
werden.

\vspace{3em}
    \captionof{figure}{Ein Rettungskorsett.
    \gSourcePicture{Wikipedia, "`Onthost"', Public domain}}
    }{
    \begin{center}
    \includegraphics[width=5cm]{./images/geraete/KED.jpg}
   

    \end{center}
}


   
\clearpage
\section{Schienungsmaterial}

\subsection{HWS-Schiene}
\index{HWS-Schiene!Beschreibung}

\gMarried{
Dient zum Ruhigstellen der {\bfseries H}als{\bfseries
W}irbel{\bfseries S}äule (HWS). Es sind verschiedene
Fabrikate im Einsatz (z.B. Stifneck\texttrademark). Es gibt
Schienen mit einer fixen Größe oder moderne Schienen, bei
welchen sich die gewümschte Größe einstellen lässt
(Stifneck select{\texttrademark} bzw. Stifneck Paedi
select{\texttrademark} für Kinder).

\vspace{3em}

    \captionof{figure}{\gAuthNotesPicLegalNotOk{GER-stiffneck.jpg}{}HWS-Schiene} 
    }{
    \begin{center}
    \includegraphics[width=0.7\linewidth]{./images/noauth/GER-stiffneck.jpg}
    \end{center}
}



\subsection{SAM-Splint}

\gMarried{
Dient am idealsten für das Ruhigstellen von Extremitätenverletzungen mit abnormer Stellung. Da es sich leicht an den
Arm an formen lässt. Es muss jedoch noch mit einer Mullbinde
oder mit Dreiecktücher fixiert werden.
}{
    \begin{center}   
    \includegraphics[width=0.7\linewidth]{./images/noauth/GER-samsplint.png}

    \captionof{figure}{\gAuthNotesPicLegalNotOk{GER-samsplint.png}{}Sam-Splint}
        \end{center}
}


\subsection{Vakuumbeinschiene}

\gMarried{
Dient zum Ruhigstellen von Verletzungen zwischen Knie und Zehen.

\vspace{3em}

    \captionof{figure}{\gAuthNotesPicLegalOk{}{}Vakuumbeinschiene.}
}{
    \begin{center}
\includegraphics[width=\linewidth]{./images/geraete/vakuumbeinschiene.jpg}
    \end{center}
}


\subsection{Vakuummatratze}

\gMarried{
Wird zum Ruhigstellen von Wirbelsäule, Becken und Oberschenkel verwendet.

\vspace{3em}

    \captionof{figure}{Vakuummatratze.}
} 
{
    \begin{center}
    \includegraphics[width=5cm]{./images/noauth/GER-vakuummatratze.png}
    \end{center}
}


% \subsubsection{Pulsoxymetrie}
% 
% \index{Pulsoxymetrie}\begin{center}
%  [Warning: Image ignored] % Unhandled or unsupported graphics:
% %\includegraphics[width=3.516cm,height=3.516cm]{AASdev20090228-img13.png}
% \gTempPicMissing{Pulsoxy img13}
% 
% \end{center}
% Das Pulsoxymeter (kurz: Pulsoxy) mi{\ss}t dden Puls und die
% Sauerstoffsättigung (\index{SpO2}SpO2)des Blutes (Verhältnis
% O2-beladene/ nichtbeladene Erythrozyten). 
% 
% Normalwert 96-99\%, 
% 
% [F078?] Achtung: Ein {\quotedblbase}guter{\textquotedblleft}
% \index{SpO2}Sp$O_2$-Wert bedeutet nicht {\quotedblbase}keinen Sauerstoff
% geben{\textquotedblleft}. Die Entscheidung bezüglich Sauerstoffgabe
% berücksichtigt die Diagnose und den Zustand des Patienten, der
% \index{SpO2}Sp$O_2$-Wert hat wenig bis gar nichts zu sagen! 
% 
% Achtung: Das Pulsoxymeter verwechselt $O_2$ und
% $CO$-Moleküle, bei einer $CO$-Vergiftung zeigt das Pulsoxy
% falsch-hohe Sättigungswerte an!
% 
% Weitere Fehlerquellen: kein Signal: Nagellack, Zentralisation
% 
% Das Pulsoxymeter misst die Herzfrequenz und die Sauerstoffsättigung
% (=Sp$O_2$) des Blutes, aber es gilt:
% 
% [F078?] Ein {\quotedblbase}guter{\textquotedblleft} Sättigungswert ist
% keine Begründung keinen Sauerstoff zu verabreichen!
\section{Das EKG}

Das Elektrokardiogramm (EKG) zeigt mittels eines entsprechenden
Gerätes die elektrische Herzaktivität des
Reizbildungs- und Reizleitungssystems an. Dazu werden auf
dem K\"orper nach einem bestimmten Muster Elektroden
aufgeklebt und mittels eines Kabels mit dem EKG-Gerät
verbunden. Im Rettungsdienst werden fast ausschließlich
Klebeelektroden für den einmaligen Gebrauch verwendet. 

\gAuthNotes{GABS: Sollte im Defi-Teil durchgenommen
werden?:} \gZ{EKG -- Ablauf einer Erregung

\begin{itemize}
    \item P-Welle \~{} Vorhoferregung
    \item PQ-Zeit \~{} AV-\"Uberleitung
    \item QRS-Komplex \~{} Kammererregung
    \item ST-Strecke + T-Welle \~{} Erregungsrückbildung
\end{itemize}}

\gCave{Das EKG gibt nur die elektrische Leitung, nicht aber
die tatsächliche Muskelarbeit wieder! \gEs{Es erlaubt
keine Aussage über die Herzleistung! }}

\gFazit{Grundsatz: \gEs{Selbst ein toter Patient kann ein
{\quotedblbase}sch\"ones{\textquotedblleft} EKG haben.}
(allerdings nicht lange)}

\subsection{EKG -- Ableitungen}
Grundsätzlich unterscheidet man zwischen den Extremitätenableitungen
und den Brustwandableitungen. 

\begin{itemize}
    \item Extremitätenableitungen: I, II, III, aVL, aVR, aVF
    \item Brustwandableitungen: V1 -- V6
\end{itemize}

Es haben sich folgende Begriffe eingebürgert:

\begin{description}
    \item[Extremitätenableitung] "`\gT{Rhythmusstreifen}"'
    zur Beurteilung des Herzrhythmus und zum Monitoring
    \item[Extremitätenableitung\,+\,Brustwandableitung]
    "`\gT{12-Kanal-EKG}"' -- Damit können Schäden der
    jeweiligen Herzregion zugeordnet werden. Wenn vorhanden,
    ist das 12-Kanal-EKG (\emph{"`12er"'}) für die
    Diagnostik das Mittel der Wahl. 
\end{description}

\subsubsection{Extremitätenableitungen}

Die Extremitätenableitungen werden mittels 4 Elektroden abgeleitet. Je
eine Elektrobe wird an den Enden der Extremitäten angebracht. Zur
Vereinfachung k\"onnen die Elektroden auch am \"Ubergang vom Rumpf zu
den Extremitäten angebracht werden (z.B. statt dem rechten Arm an der
rechten Schulter). Wichtig ist, dass die 
Elektrodenpärchen immer auf \gEs{gleicher Höhe} angebracht
werden!

Aus den Informationen der 4 Elektroden erzeugt das
EKG-Gerät \gZ{6} Ableitungen\gZ{:  I, II, III, aVL, aVR
und aVF.}

\begin{center}
\includegraphics[width=7cm]{./images/ECGcolor.png}
\captionof{figure}{\gAuthNotesPicLegalOk{}{PD}Extremitätenableitungen
(schwarze Kabel) und Brustwandableitungen (blaue Kabel).
\gSourcePicture{Madhero88 (Wikimedia Commons), Public
domain}\gAuthNotes{GABS: Bessere Beschriftung wäre
wünschenswert, Dringlichkeit: niedrig}}
\end{center}


 
\gAuthNotes{Elektroden noch farblich Kennzeichnen}
\begin{center}
\begin{tabulary}{\linewidth}{CCC}\toprule
\gTabHeading{Elektrode} & \gTabHeading{Position} &
\gTabHeading{Alternative Postion}
\\\midrule
\gTabSub{\color{red}Rot} & Rechte Hand & Rechte
Schulter
\\
\gTabSub{\color{yellow}Gelb} & Linke Hand & Linke
Schulter
\\
\gTabSub{\color{darkGreen}Grün } & Linker Fu{\ss}  & Linke
Leiste
\\
\gTabSub{\color{black}Schwarz} & Rechter Fu{\ss} & Rechte
Leiste \\\bottomrule
\end{tabulary}
\end{center}

\subsubsection{Brustwandableitungen}

Die Brustwandableitungen werden mittels 6 Elektroden abgeleitet und
erlauben eine genauere Zuordnung von St\"orungen zu der jeweiligen
Region des Herzens. Aus den Informationen der 6 Elektroden erzeugt das
EKG-Gerät 6 Ableitungen: V1 -- V6. 

\begin{center}
\begin{tabulary}{\linewidth}[H]{cL}
\toprule
Elektrode
 &
Position
\\\midrule
V1 & 4. Zwischenrippenraum rechts vom Sternum
\\
V2 & 4. Zwischenrippenraum links vom Sternum
\\
V3 & Mitte zwischen V2 und V4
\\
V4 & 5. Zwischenrippenraum links in der Verlängerung der
Schlüsselbein-Mittellinie
\\
V5 & Mitte zwischen V4 und V6
\\
V6 & 6. Zwischenrippenraum in der mittleren Achsellinie
\\\bottomrule
\end{tabulary}
\end{center}


\section{Pulsoxymetrie}

Die Pulsoxymetrie dient zur nicht-invasiven Überwachung der
Sauerstoffsättigung (Sp$O_2$) des arteriellen Blutes. Dies
geschieht mittels Photorezeptoren und Infrarotlicht. Dabei
wird das Verhältnis der mit Sauerstoff beladenen roten
Blutkörperchen zu den nicht beladenen bestimmt. Der
Normalwert liegt zwischen 96\% und 99\%. Als zweiten Wert
zeigt das Pulsoxymeter die aktuelle Pulsfrequenz an.

Mögliche Fehlerquellen, wenn das Gerät keinen brauchbaren Wert anzeigt:
\begin{itemize}
	\item CO-Vergiftungen; Das Gerät verwechselt $O_2$ und
	$CO$-Moleküle und zeigt einen falschen hohen Wert.
	\item Minderdurchblutung der Peripherie; Entweder durch Zentralisation beim Schock oder Unterkühlung.
	\item Nagellack
	\item künstliche Fingernägel
	\item Sonneneinstrahlung von der Seite
\end{itemize}

\gFazit{Immer zuerst auf den Patienten schauen (Symptome)
und erst dann auf das Pulsoxymeter!}


